% from ega2-1-fr.tex
\oldpage[II]{95}

  \item[\rm{e)}] Tout $\mathcal{O}_X$-Module quasi-cohérent est engendré par
  ses sections au-dessus de $X$.
  \item[\rm{e')}] Tout faisceau quasi-cohérent d'idéaux de type fini de
  $\mathcal{O}_X$ est engendré par ses sections au-dessus de $X$.
\end{enumerate}
\end{proposition}

Il est clair que b) entraîne a), et a) entraîne c) en vertu du critère b) de
(\hyperref[II.4.4.4-fr]{4.4.4}) appliqué au morphisme identique (compte tenu
de la remarque précédant l'énoncé)~; d'autre part c) entraîne c')
(\hyperref[II.4.5.10-fr]{4.5.10, (i)}), et c'), b) et b') sont équivalents par
(\hyperref[II.4.5.2-fr]{4.5.2, critères b) et b')}). Enfin c') équivaut à
chacun des critères d), d'), e), e') en vertu de
(\hyperref[II.4.5.2-fr]{4.5.2, critères a), a'), c)}) et
(\hyperref[II.4.5.5-fr]{4.5.5, critère d'')}).

On observera en outre qu'avec les notations précédentes, ceux des $X_f$ qui
sont affines forment une \emph{base} de la topologie de $X$, et que le
morphisme canonique $u$ est \emph{dominant}
(\hyperref[II.4.5.2-fr]{4.5.2}).

\begin{corollary}[5.1.3]
\label{II.5.1.3-fr}
Soit $X$ un préschéma quasi-compact. S'il existe un morphisme $v:X\to Y$ de
$X$ dans un schéma affine $Y$, qui soit un homéomorphisme de $X$ sur un
sous-espace ouvert de $Y$, alors $X$ est quasi-affine.
\end{corollary}

En effet, il existe une famille $(g_\alpha)$ de sections de $\mathcal{O}_Y$
au-dessus de $Y$ telles que les $D(g_\alpha)$ forment une base de la
topologie de $v(X)$~; si $v=(\psi,\theta)$ et si l'on pose
$f_\alpha=\theta(g_\alpha)$, on a
$X_{f_\alpha}=\psi^{-1}(D(g_\alpha))$
(\hyperref[I.2.2.4.1-fr]{I, 2.2.4.1}), donc les $X_{f_\alpha}$ forment une
base de la topologie de $X$, et le critère d) de
(\hyperref[II.5.1.2-fr]{5.1.2}) est vérifié.

\begin{corollary}[5.1.4]
\label{II.5.1.4-fr}
Si $X$ est un schéma quasi-affine, tout $\mathcal{O}_X$-Module inversible est
très ample (relativement au morphisme canonique) et a fortiori ample.
\end{corollary}

En effet, un tel Module $\mathcal{L}$ est engendré par ses sections au-dessus
de $X$ (\hyperref[II.5.1.2-fr]{5.1.2, e)}), donc
$\mathcal{L}\otimes\mathcal{O}_X=\mathcal{L}$ est très ample
(\hyperref[II.4.4.8-fr]{4.4.8}).

\begin{corollary}[5.1.5]
\label{II.5.1.5-fr}
Soit $X$ un préschéma quasi-compact. S'il existe un $\mathcal{O}_X$-Module
inversible $\mathcal{L}$ tel que $\mathcal{L}$ et $\mathcal{L}^{-1}$ soient
amples, $X$ est un schéma quasi-affine.
\end{corollary}

En effet, $\mathcal{O}_X=\mathcal{L}\otimes\mathcal{L}^{-1}$ est alors ample
(\hyperref[II.4.5.7-fr]{4.5.7}).

\begin{proposition}[5.1.6]
\label{II.5.1.6-fr}
Soit $f:X\to Y$ un morphisme quasi-compact. Les conditions suivantes sont
équivalentes~:
\begin{enumerate}
  \item[\rm{a)}] $f$ est quasi-affine.
  \item[\rm{b)}] La $\mathcal{O}_Y$-Algèbre
  $f_*(\mathcal{O}_X)=\mathcal{A}(X)$ est quasi-cohérente et le morphisme
  canonique $X\to\operatorname{Spec}(\mathcal{A}(X))$ correspondant à
  l'homomorphisme identique $\mathcal{A}(X)\to\mathcal{A}(X)$
  (\hyperref[II.1.2.7-fr]{1.2.7}) est une immersion ouverte.
  \item[\rm{b')}] La $\mathcal{O}_Y$-Algèbre $\mathcal{A}(X)$ est
  quasi-cohérente et le morphisme canonique
  $X\to\operatorname{Spec}(\mathcal{A}(X))$ est un homéomorphisme de $X$
  sur un sous-espace de $\operatorname{Spec}(\mathcal{A}(X))$.
  \item[\rm{c)}] Le $\mathcal{O}_X$-Module $\mathcal{O}_X$ est très ample
  pour $f$.
  \item[\rm{c')}] Le $\mathcal{O}_X$-Module $\mathcal{O}_X$ est ample pour
  $f$.
  \item[\rm{d)}] Le morphisme $f$ est séparé et, pour tout
  $\mathcal{O}_X$-Module quasi-cohérent $\mathcal{F}$, l'homomorphisme
  canonique $\sigma:f^*(f_*(\mathcal{F}))\to\mathcal{F}$
  (\hyperref[0.4.4.3-fr]{0, 4.4.3}) est surjectif.
\end{enumerate}
En outre, lorsque $f$ est quasi-affine, tout $\mathcal{O}_X$-Module
inversible $\mathcal{L}$ est très ample relativement à $f$.
\end{proposition}

L'équivalence de a) et c') résulte du caractère local sur $Y$ de la
$f$-amplitude (\hyperref[II.4.6.4-fr]{4.6.4}), de la définition
(\hyperref[II.5.1.1-fr]{5.1.1}) et du critère
(\hyperref[II.5.1.2-fr]{5.1.2, c')}). Les autres propriétés sont locales sur
$Y$

